% Section 3: The quadratic trichotomy and the central sign
\label{sec:scalar-body}

A recurring temptation in this material is to declare the complex field as the
foundational scalar and then to derive the geometry.  This section shows why that
order is unavailable.  Within the natural one-parameter family of two-dimensional real
algebras, the coarse requirements that a "characteristic" theory might impose are
satisfied by all three members; only a genuinely geometric condition separates them,
and that condition has a precise meaning: it is the retention of a central sign.

\subsection{The quadratic family}

\begin{definition}[Quadratic scalar family]\label{def:quadratic-family}
For $\kappa\in\{-1,0,1\}$ let
\[
  A_\kappa=\R[j]/(j^{2}-\kappa),
\]
so that $a+bj$ with $a,b\in\R$ multiplies by
\[
  (a+bj)(c+dj)=(ac+\kappa bd)+(ad+bc)j .
\]
\end{definition}

\begin{proposition}[Elementary structure]\label{prop:quadratic-structure}
In each $A_\kappa$: the quadratic norm $N_\kappa(x)=x\bar x$ with
$\overline{a+bj}=a-bj$ is multiplicative, and conjugation is a multiplicative
involution.  Moreover for every $u\in A_\kappa$ the rank-one map $L_u(x)=ux$
satisfies
\[
  L_u(1)=u,\qquad L_u\circ L_v=L_{uv},
\]
so a scalar is recovered from its own characteristic action without coordinates, and
the three aspect units $\tau,\sigma,\rho\in A_\kappa$ compose cyclically to a single
characteristic $L_{\rho\sigma\tau}$.
\end{proposition}

\begin{proof}
Multiplicativity of $N_\kappa$ is the identity
$(ac+\kappa bd)^{2}-\kappa(ad+bc)^{2}=(a^{2}-\kappa b^{2})(c^{2}-\kappa d^{2})$,
which is a direct expansion.  The rest is immediate from associativity of $A_\kappa$
and $L_u(1)=u$.
\end{proof}

The three cases are conventionally named by the sign of the norm:

\begin{center}
\begin{tabular}{@{}ccll@{}}
\toprule
$\kappa$ & algebra & $N_\kappa(a,b)$ & type \\
\midrule
$-1$ & complex numbers & $a^{2}+b^{2}$ & elliptic \\
$0$ & dual numbers & $a^{2}$ & parabolic \\
$1$ & split-complex numbers & $a^{2}-b^{2}$ & hyperbolic \\
\bottomrule
\end{tabular}
\end{center}

\begin{remark}[What Proposition~\ref{prop:quadratic-structure} does not do]
\label{rem:not-enough}
It does not select $\kappa=-1$.  Rank-one recovery of a characteristic from its own
action, multiplicative conjugation, and a multiplicative quadratic norm are all
satisfied by the three members.  Any argument that begins "the scalar must be
complex, since only then can one conjugate and take norms" is therefore invalid.
\end{remark}

\subsection{The three nondegeneracy conditions, and why they are not derivations}

Three further requirements do separate the cases.

\begin{proposition}\label{prop:three-separators}
Within $\{A_\kappa\}$ each of the following conditions holds for $\kappa=-1$ alone:
\begin{enumerate}
\item every nonzero element is invertible;
\item $N_\kappa$ is positive definite;
\item the unit-norm group $\{x:N_\kappa(x)=1\}$ is compact.
\end{enumerate}
\end{proposition}

\begin{proof}
For $\kappa=-1$ the inverse is $(a+bj)^{-1}=(a-bj)/(a^{2}+b^{2})$ and the unit group
is $S^{1}$.  For $\kappa=0$, every $bj$ with $b\neq0$ is a noninvertible nonzero
element.  For $\kappa=1$, $(1+j)(1-j)=0$ exhibits nonzero zero divisors and the norm
is indefinite with noncompact unit locus.
\end{proof}

\begin{warning}
Proposition~\ref{prop:three-separators} is a separation result, not a derivation.
None of the three conditions is forced by the geometry of readings.  They are
\emph{available} selection principles, and choosing one is a mathematical decision
that must be declared.
\end{warning}

\subsection{The negative-inverse criterion}

A better candidate comes from the geometry itself.  Paper~0 already contains the
M\"obius involution $S(z)=-1/z$ as the reciprocal reading
\cite{Yuan2026AEGPaperZero}; the question is which quadratic scalars admit it
globally.

\begin{proposition}[Negative inverse in the quadratic family]\label{prop:negative-inverse}
For $z=a+bj\in A_\kappa$ with $a^{2}-\kappa b^{2}\neq0$,
\[
  S_\kappa(a,b)=\Bigl(\frac{-a}{a^{2}-\kappa b^{2}},\ \frac{b}{a^{2}-\kappa b^{2}}\Bigr)
\]
satisfies $z\,S_\kappa(z)=-1$.  For $\kappa=-1$ the formula
\begin{enumerate}
\item is defined on the whole punctured plane $\R^{2}\setminus\{0\}$;
\item is an involution;
\item satisfies $zS_{-1}(z)=-1$ identically;
\item preserves the upper half-plane $\{b>0\}$.
\end{enumerate}
For $\kappa=0$ the entire nilpotent axis $a=0$ is excluded.  For $\kappa=1$ the two
light rays $a=\pm b$ are excluded, and the sign of the second coordinate need not be
preserved across the components they separate.
\end{proposition}

\begin{proof}
$1/z=(a-bj)/(a^{2}-\kappa b^{2})$, so $-1/z=S_\kappa(a,b)$ as displayed, and
$zS_\kappa(z)=-1$ follows.  For $\kappa=-1$ the denominator $a^{2}+b^{2}$ never
vanishes off the origin, and $b/(a^{2}+b^{2})>0$ whenever $b>0$; involutivity is the
usual $(S\circ S)(z)=z$ for a M\"obius involution on a chart where it is defined.  The
excluded loci for $\kappa=0,1$ are exactly the zero divisors.
\end{proof}

\begin{proposal}[The AEG inversion criterion]\label{prop:inversion-criterion}
Among the tested family, $\C$ is selected by the requirement that a globally defined
negative-inverse involution exist on every nonzero scalar and preserve one connected
real half-plane.
\end{proposal}

This is labelled \textsf{Proposal} because the requirement is imported from the
observed transition between two real grids, not derived from program geometry; what
is proved is only that it selects.

\subsection{The oriented lift, and chirality as a central sign}

The inversion criterion has a sharper form which explains \emph{why} it selects, and
which supplies the paper's working definition of chirality.

\begin{definition}[Time--space exchange and its lift]\label{def:exchange-lift}
On a homogeneous pair $(t,s)$ put
\[
  J(t,s)=(-s,t),\qquad J=\begin{pmatrix}0&-1\\ 1&0\end{pmatrix}.
\]
On the projective chart $z=t/s$ this induces $[t:s]\mapsto[-s:t]$, i.e.
$z\mapsto-1/z$.  We call $J$ the \emph{oriented lift} of the projective exchange.
\end{definition}

\begin{proposition}[Order two versus order four]\label{prop:order-two-four}
$[J]^{2}=[I]$ on the projective line, while $J^{2}=-I$ on the oriented lift.  Hence
$J^{-1}=-J$, and the forward and reverse exchanges induce the same projective
involution although they are distinct on the lift.
\end{proposition}

\begin{proof}
$J^{2}=\begin{pmatrix}-1&0\\0&-1\end{pmatrix}$.  Projectively $[-\id]=[\id]$, so
$[J]^{2}=[\id]$, while on the lift the scalar $-1$ is retained.
\end{proof}

\begin{definition}[Chirality]\label{def:chirality}
\emph{Chirality} is the central sign distinguishing the forward and reverse oriented
lifts of one projective exchange: the datum that $J$ and $J^{-1}=-J$ are different
lifts of the same projective involution.
\end{definition}

This definition is deliberately austere.  It says nothing about the orientation of the
ambient plane, and the following computation shows that it must not.

\begin{proposition}[Chirality is not ambient orientation reversal]\label{prop:not-orientation}
In real coordinates $S(x,y)=\bigl(-x/(x^{2}+y^{2}),\,y/(x^{2}+y^{2})\bigr)$ has
Jacobian determinant
\[
  \det DS=\frac{1}{(x^{2}+y^{2})^{2}}>0 .
\]
Thus the involution is holomorphic and orientation-preserving on the plane; the
chirality change is not a reversal of that plane but the difference between the two
lifts of one projective map.
\end{proposition}

\begin{proof}
$S$ is the M\"obius map $-1/z$ with derivative $1/z^{2}$, so
$\det DS=|z|^{-4}=(x^{2}+y^{2})^{-2}>0$.
\end{proof}

\subsection{The generator action under one bridge hypothesis}

The central sign becomes a selection mechanism for the algebra when the same oriented
two-dimensional carrier is asked to represent both the exchange lift and the action of
the quadratic generator.

\begin{proposition}[Generator action]\label{prop:generator-action}
Let $L_j$ be multiplication by $j$ on $(a,b)$, so that
\[
  L_j=\begin{pmatrix}0&\kappa\\ 1&0\end{pmatrix},\qquad L_j^{2}=\kappa I .
\]
Then:
\begin{enumerate}
\item for $\kappa=-1$, $L_j=J$ is orientation-preserving, invertible, of order four
  before projectivization;
\item for $\kappa=1$, $L_j^{2}=+I$, $\det L_j=-1$, and the generator is already its own
  inverse on the lift;
\item for $\kappa=0$, $L_j$ is singular and nilpotent.
\end{enumerate}
\end{proposition}

\begin{proof}
Direct computation for each $\kappa$; $\det L_j=-\kappa$ and $L_j^{2}=\kappa I$.
\end{proof}

\begin{calibration}[The bridge is a hypothesis]\label{cal:bridge-hypothesis}
Requiring an oriented time--space exchange whose projective class has order two while
its lift distinguishes forward from reverse selects $j^{2}=-1$, \emph{provided} the
same oriented two-dimensional carrier is declared to represent both the exchange lift
and the generator action.  Both halves are calibrated, not proved: the selection is a
finite algebraic check, and the bridge is a declared identification.
\end{calibration}

\begin{warning}
The matrix $J\in SL(2,\R)$ already defines the real-coefficient M\"obius map $-1/z$;
complex notation is therefore not by itself evidence that the underlying scalars are
complex.  What is new is that retaining the lift rather than quotienting to
$\PSL(2,\R)$ exposes the central sign.  Calling that sign chirality requires the
oriented carrier to remember it.
\end{warning}

\subsection{The conjugated presentation}

The last piece of the trichotomy is what the involution does to the two AEG
generators, because it shows the central sign in operational form.

\begin{proposition}[Conjugated generators]\label{prop:conjugated-generators}
Write $A(z)=z+1$ and $M_\lambda(z)=\lambda z$ for $\lambda>0$ in the upper
half-plane chart.  Then
\[
  S M_\lambda S=M_{\lambda^{-1}},\qquad
  S A S(z)=\frac{z}{1-z}.
\]
Thus conjugation inverts the scale direction while turning the translation grid into a
parabolic grid organised around the other fixed point.  For a Baumslag--Solitar-type
relation $MAM^{-1}=A^{n}$ the conjugated presentation has scale generator $M^{-1}$ and
ripple generator $SAS$.
\end{proposition}

\begin{proof}
For the first identity, $S M_\lambda S(z)=-1/(\lambda\cdot(-1/z))=z/\lambda
=M_{\lambda^{-1}}(z)$.  For the second,
$SAS(z)=-1/((-1/z)+1)=z/(1-z)$, which is parabolic with the single fixed point $1/2$.
\end{proof}

\begin{remark}[Baumslag--Solitar background]
The relation $MAM^{-1}=A^{n}$ is the Baumslag--Solitar presentation
\cite{BaumslagSolitar1962NonHopfian}; Paper~0 records the exact routing of
copy/branch/recombine programs through these relations.  The content here is only that
conjugation exchanges the forward and reverse presentations of the selected scale
direction, which is the visible consequence of lift-level chirality.
\end{remark}

\subsection{What the trichotomy settles, and what it leaves open}

\begin{calibration}[Bounded red-team calibration]\label{cal:trichotomy-bounded}
Within the tested quadratic family: two real dimensions per aspect, cyclic rank-one
action, recovery of a characteristic from its unit action, conjugation, and a
multiplicative norm do \emph{not} select $\C$; a global negative-inverse involution
preserving a connected real AEG half-plane \emph{does} select $\C$; and the
mechanism is that the projective exchange has order two while its
orientation-preserving lift has order four.
\end{calibration}

Three interpretations of the six real degrees of freedom remain live and are not
separated by anything proved here.

\begin{proposal}[Three readings of $6=2+2+2$]\label{prop:three-readings}
The cases $\kappa=-1,0,1$ may respectively encode phase/orientation, infinitesimal
extension, and causal scale; alternatively the three pairs may be three complex
lines; alternatively $\C$, dual, and split behaviour may be charts or strata of one
expression-valued scalar geometry.  No canonical assignment of the three cases to
temporal, spatial, and relational roles is established.
\end{proposal}

\begin{remark}[The Jordan shoulder]
The dual case is not a pathology to be eliminated.  A nilpotent extension naturally
records a first-order shoulder that survives complexification, which suggests that a
faithful spectrum may require both a semisimple character and a residual extension
object.  \S\ref{sec:residual-catalogue} returns to this.
\end{remark}
